\begin{tikzpicture}[font=\small,line cap=round,line join=round]\geoEnonce{\node[font=\tiny,otMuted]at(2.4,1.0){écrire le programme, puis l’exécuter en vraie grandeur};}\geoSolution{\begin{scope}[shift={(2.4,0)}]\draw[geoReponse](0,0)circle(1.5);\draw[geoReponse](1.5,0.0)--(0.7500000000000002,1.299038105676658)--(-0.7499999999999997,1.299038105676658)--(-1.5,1.8369701987210297e-16)--(-0.7500000000000007,-1.2990381056766578)--(0.7500000000000002,-1.299038105676658)--cycle;\node[below,font=\tiny,otGreenDark]at(0,-1.85){cercle de 3 cm, rayon reporté six fois};\end{scope}}\end{tikzpicture}\par\vspace{1.5mm}{\otRegular\fontsize{8}{10}\selectfont\color{otMuted}Le programme : un cercle, six reports du rayon, six segments. Échelle réduite au corrigé.\par}
